% English translation of chapter1.tex lines 1708--2079.
% Source: Methods of Algebra, Volume 2, commit
% 9a5803ff77dd3257484cb177f851a73770a59dd3 (CC BY 4.0).
% stable-unit-id: o014.aljabr2.chapter1.adjoint-functor-theorem

% segment-id: o014.li2.chapter1.adjoint-functor-theorem.h001
\section{Adjoint Functor Theorems}\label{sec:AFT}
\hypertarget{o014.aljabr2.chapter1.adjoint-functor-theorem}{ }

% segment-id: o014.li2.chapter1.adjoint-functor-theorem.p001
Functors often occur in adjoint pairs. Consider categories $\mathcal{C}$ and
$\mathcal{D}$ and a pair of functors
% segment-id: o014.li2.chapter1.adjoint-functor-theorem.g001
\[\begin{tikzcd}
	F: \mathcal{C} \arrow[shift left, r] & \mathcal{D}: G \arrow[shift left, l]
\end{tikzcd}\]
such that $(F,G)$ is an adjoint pair. Then $F$ preserves $\varinjlim$ and $G$
preserves $\varprojlim$, whenever the limits in question exist. Conversely,
one naturally asks: if $F:\mathcal{C}\to\mathcal{D}$ preserves colimits,
what ensures that it has a right adjoint? If
$G:\mathcal{D}\to\mathcal{C}$ preserves limits, what ensures that it has a
left adjoint? In many cases the adjoint can be written down explicitly, but
sometimes an abstract existence result is needed. Such results are
collectively called adjoint functor theorems. Their hypotheses fall roughly
into two classes. First, to use preservation of colimits (or limits), one
assumes that $\mathcal{C}$ (or $\mathcal{D}$) is cocomplete (or complete).
Second, one needs conditions closely related to set size.
\index{size issues}

% segment-id: o014.li2.chapter1.adjoint-functor-theorem.p002
Our first task is to reformulate the existence of an adjoint appropriately.
By duality, we state only the case of a left adjoint.

% segment-id: o014.li2.chapter1.adjoint-functor-theorem.p003
For a functor $G:\mathcal{D}\to\mathcal{C}$ and
$c\in\Obj(\mathcal{C})$, Definition \ref{def:comma-category} gives a
category $(c/G)$ whose objects are data $(d,c\xrightarrow{f}Gd)$.

% segment-id: o014.li2.chapter1.adjoint-functor-theorem.le001
\begin{lemma}\label{prop:adjoint-vs-comma}
	A functor $G:\mathcal{D}\to\mathcal{C}$ has a left adjoint if and only if
	$(c/G)$ has an initial object for every $c\in\Obj(\mathcal{C})$.
\end{lemma}
% segment-id: o014.li2.chapter1.adjoint-functor-theorem.pf001
\begin{proof}
	Fix $c$. Specifying a morphism
	$(d_0,c\xrightarrow{\beta_0}Gd_0)\to(d,c\xrightarrow{\beta}Gd)$ in
	$(c/G)$ amounts to specifying $\alpha:d_0\to d$ such that
% segment-id: o014.li2.chapter1.adjoint-functor-theorem.g002
	\[\begin{tikzcd}[row sep=small, column sep=small]
		Gd_0 \arrow[rr, "G\alpha"] & & Gd \\
		& c \arrow[lu, "\beta_0"] \arrow[ru, "\beta"'] &
	\end{tikzcd}\]
	commutes. Thus $(d_0,\beta_0)$ is initial in $(c/G)$ if and only if
% segment-id: o014.li2.chapter1.adjoint-functor-theorem.q001
	\begin{equation}\label{eqn:adjoint-vs-comma-aux}
% segment-id: o014.li2.chapter1.adjoint-functor-theorem.g003
	\begin{tikzcd}[row sep=tiny]
		\Hom_{\mathcal{D}}(d_0, d) \arrow[r] & \Hom_{\mathcal{C}}(c, Gd) \\
		\alpha \arrow[mapsto, r] & \beta := (G\alpha) \beta_0
	\end{tikzcd}\end{equation}
	is a bijection for every $d\in\Obj(\mathcal{D})$. This is the basis of the
	argument.

	Suppose that $G$ has a left adjoint $F$, and let
	$c\in\Obj(\mathcal{C})$. We claim that $Fc$, together with the unit
	$\eta_c:c\to G(Fc)$, is initial in $(c/G)$. For any
	$(d,c\xrightarrow{\beta}Gd)$, the map in
	\eqref{eqn:adjoint-vs-comma-aux},
	$\Hom_{\mathcal{D}}(Fc,d)\to\Hom_{\mathcal{C}}(c,Gd)$, sends $\alpha$ to
	$(G\alpha)\eta_c$. This is exactly the bijection on $\Hom$ sets supplied
	by the adjunction, so $(Fc,\eta_c)$ is initial.

	Conversely, suppose that $(c/G)$ has an initial object for every $c$, and
	write it as $(Fc,\eta_c:c\to G(Fc))$. Since initial objects are unique up
	to unique isomorphism, these objects form a functor
	$F:\mathcal{C}\to\mathcal{D}$, and
	$(\eta_c)_{c\in\Obj(\mathcal{C})}$ gives
	$\eta:\identity_{\mathcal{C}}\to GF$. For every
	$(c,d)\in\Obj(\mathcal{C})\times\Obj(\mathcal{D})$, define the bijection
	from \eqref{eqn:adjoint-vs-comma-aux}:
% segment-id: o014.li2.chapter1.adjoint-functor-theorem.q002
	\begin{align*}
		\varphi_{c, d}: \Hom_{\mathcal{D}}(Fc, d) & \to \Hom_{\mathcal{C}}(c, Gd) \\
		\alpha & \mapsto (G\alpha) \eta_c.
	\end{align*}
	These bijections are plainly functorial in $c$ and $d$. Hence
	$(F,G,\varphi)$ is an adjoint pair with unit $\eta$.
\end{proof}

% segment-id: o014.li2.chapter1.adjoint-functor-theorem.d001
\begin{definition}
	Let $\mathcal{E}$ be a category. A nonempty small subset $\Gamma$ of
	$\Obj(\mathcal{E})$ is called a \emph{weakly initial family} if, for every
	$e'\in\Obj(\mathcal{E})$, there is an $e\in\Gamma$ such that
	$\Hom_{\mathcal{E}}(e,e')$ is nonempty.\footnote{Translator's note: the
	source has the subscript $\mathcal{C}$; it has been corrected to
	$\mathcal{E}$, the category introduced in this definition (O014-C020).}
	If $e\in\Obj(\mathcal{E})$ and $\{e\}$ is weakly initial, then $e$ itself
	is called a \emph{weakly initial object}.
\end{definition}

% segment-id: o014.li2.chapter1.adjoint-functor-theorem.p004
Every initial object is weakly initial. Weakly terminal families and weakly
terminal objects are defined similarly.

% segment-id: o014.li2.chapter1.adjoint-functor-theorem.le002
\begin{lemma}\label{prop:weakly-initial-prod}
	If $\Gamma$ is a weakly initial family in $\mathcal{E}$ and
	$e:=\prod_{x\in\Gamma}x$ exists, then $e$ is weakly initial.
\end{lemma}
% segment-id: o014.li2.chapter1.adjoint-functor-theorem.pf002
\begin{proof}
	For every $e'\in\Obj(\mathcal{E})$, there are $y\in\Gamma$ and a morphism
	$y\to e'$. Take the composite
	$e\xrightarrow{\text{projection}}y\to e'$.
\end{proof}

% segment-id: o014.li2.chapter1.adjoint-functor-theorem.p005
Return to adjoint functor theorems. To apply Lemma
\ref{prop:adjoint-vs-comma}, we must study the projection functor
$\Pi_{c/}:(c/G)\to\mathcal{D}$ from Definition \ref{def:comma-category}.

% segment-id: o014.li2.chapter1.adjoint-functor-theorem.pr001
\begin{proposition}\label{prop:comma-cat-completeness}
	Let $\mathcal{D}$ be complete, let $G:\mathcal{D}\to\mathcal{C}$ preserve
	all small limits, and fix $c\in\Obj(\mathcal{C})$.
% segment-id: o014.li2.chapter1.adjoint-functor-theorem.l001
	\begin{enumerate}[(i)]
% segment-id: o014.li2.chapter1.adjoint-functor-theorem.i001
		\item The functor $\Pi_{c/}$ is conservative (Definition
		\ref{def:conservative}) and creates all small limits (Definition
		\ref{def:create-limit}).
% segment-id: o014.li2.chapter1.adjoint-functor-theorem.i002
		\item The category $(c/G)$ is complete.
% segment-id: o014.li2.chapter1.adjoint-functor-theorem.i003
		\item The functor $\Pi_{c/}$ sends monomorphisms to monomorphisms.
	\end{enumerate}
\end{proposition}
% segment-id: o014.li2.chapter1.adjoint-functor-theorem.pf003
\begin{proof}
	For (i), conservativity is clear. The key point is that, given
	$\beta:J\to(c/G)$ sending $j\in\Obj(J)$ to
	$(d_j,c\xrightarrow{f_j}Gd_j)$, one can lift $\varprojlim_jd_j$ and its
	projection morphisms to
	$(\varprojlim_jd_j,c\xrightarrow{f}G(\varprojlim_jd_j))$ and show that it
	is $\varprojlim\beta$. Explicitly, take $f$ to be the composite in the
	first row of the commutative diagram
% segment-id: o014.li2.chapter1.adjoint-functor-theorem.g004
	\[\begin{tikzcd}
		c \arrow[r, "{\exists!}"] \arrow[rd, "{f_{j'}}"'] & \varprojlim_j Gd_j \arrow[d] & G (\varprojlim_j d_j) \arrow[l, "\exists!"', "\sim"] \arrow[ld] \\
		& Gd_{j'} &
	\end{tikzcd} \quad j' \in \Obj(J). \]
	The remaining checks are routine and are omitted; compare the examples in
	\S\ref{sec:limit-functor}.

	Statement (ii) follows from the completeness of $\mathcal{D}$ and (i).

	For (iii), first observe that for every morphism $f:a\to b$ in any
	category, a direct manipulation of the definitions gives
% segment-id: o014.li2.chapter1.adjoint-functor-theorem.g005
	\[ f \;\text{is monic} \iff
	\begin{tikzcd}
		a \arrow[r, "\identity_a"] \arrow[d, "\identity_a"'] & a \arrow[d, "f"] \\
		a \arrow[r, "f"'] & b
	\end{tikzcd} \;\text{is a pullback diagram}. \]

	Now consider a monomorphism in $(c/G)$ from $(d',c\to Gd')$ to
	$(d,c\to Gd)$, determined by $\iota:d'\to d$ in $\mathcal{D}$. By (i),
% segment-id: o014.li2.chapter1.adjoint-functor-theorem.q003
	\[ (d', c \to Gd') \dtimes{(d, c \to Gd)} (d', c \to Gd') = \left( d' \dtimes{d} d' , \; c \to G\left( d' \dtimes{d} d' \right) \right) . \]
	The left-hand side identifies with $(d',c\to Gd')$, with both projections
	corresponding to $\identity_{d'}$. Thus $d'\dtimes{d}d'$ identifies with
	$d'$ in the same way. In other words, $\iota$ is monic.
\end{proof}

% segment-id: o014.li2.chapter1.adjoint-functor-theorem.le003
\begin{lemma}\label{prop:jointly-weak-initial-obj}
	If a complete category $\mathcal{E}$ has a weakly initial family $\Gamma$,
	then $\mathcal{E}$ has an initial object.
\end{lemma}
% segment-id: o014.li2.chapter1.adjoint-functor-theorem.pf004
\begin{proof}
	Completeness and Lemma \ref{prop:weakly-initial-prod} ensure that
	$\mathcal{E}$ has a weakly initial object $e$. Let $I$ be the full
	subcategory of $\mathcal{E}$ with $\Obj(I)=\{e\}$; it is small.
	Completeness gives a limit $x$ of the inclusion $I\to\mathcal{E}$, with
	canonical morphism $i:x\to e$. Its universal property says
% segment-id: o014.li2.chapter1.adjoint-functor-theorem.q004
	\begin{equation}\label{eqn:jointly-weak-initial-obj-aux}
% segment-id: o014.li2.chapter1.adjoint-functor-theorem.g006
	\begin{tikzcd}[row sep=tiny]
		i_*: \Hom_{\mathcal{E}}(t, x) \arrow[r, "1:1"] & \left\{ \phi \in \Hom_{\mathcal{E}}(t, e) : \forall \theta, \theta' \in \Hom_{\mathcal{E}}(e, e), \; \theta \phi = \theta' \phi \right\} \\
		\psi \arrow[r, mapsto] & i\psi
	\end{tikzcd}\end{equation}
	for arbitrary $t\in\Obj(\mathcal{E})$. This immediately shows that $i$ is
	monic.

	We show that $x$ is initial. Given $y\in\Obj(\mathcal{E})$, choose any
	morphism $e\to y$ and call the composite $x\xrightarrow{i}e\to y$ by
	$f$. For any $g:x\to y$, consider the equalizer
	$j:\Ker(f,g)\hookrightarrow x$. There is a morphism
	$k:e\to\Ker(f,g)$, hence $ijk:e\to e$. Taking $\psi=\identity_x$ in
	\eqref{eqn:jointly-weak-initial-obj-aux} gives
	$(ijk)i=(\identity_e)i=i$. Cancelling the monomorphism $i$ on the left
	gives $jki=\identity_x$, and consequently $f=fjki=gjki=g$. This proves the
	claim.
\end{proof}

% segment-id: o014.li2.chapter1.adjoint-functor-theorem.p006
The following result is due to P.\ Freyd and is also known as the General
Adjoint Functor Theorem.

% segment-id: o014.li2.chapter1.adjoint-functor-theorem.th001
\begin{theorem}[P.\ Freyd's Adjoint Functor Theorem]\label{prop:GAFT}
	\index{adjoint functor theorem}
	Let $\mathcal{C}$ and $\mathcal{D}$ be categories.
% segment-id: o014.li2.chapter1.adjoint-functor-theorem.l002
	\begin{enumerate}[(i)]
% segment-id: o014.li2.chapter1.adjoint-functor-theorem.i004
		\item Let $\mathcal{D}$ be complete and let
		$G:\mathcal{D}\to\mathcal{C}$ preserve all small limits. Then $G$ has
		a left adjoint if and only if the following \emph{solution set
		condition} holds: for every $c\in\Obj(\mathcal{C})$, there is a family
		of morphisms $(f_i:c\to Gd_i)_{i\in I}$, indexed by a small set $I$,
		such that for every $f:c\to Gd$ there are $i\in I$ and
		$\alpha:d_i\to d$ for which $f$ factors as
		$c\xrightarrow{f_i}Gd_i\xrightarrow{G\alpha}Gd$.
% segment-id: o014.li2.chapter1.adjoint-functor-theorem.i005
		\item Let $\mathcal{C}$ be cocomplete and let
		$F:\mathcal{C}\to\mathcal{D}$ preserve all small colimits. Then $F$ has
		a right adjoint if and only if the following \emph{cosolution set
		condition} holds: for every $d\in\Obj(\mathcal{D})$, there is a family
		of morphisms $(g_i:Fc_i\to d)_{i\in I}$, indexed by a small set $I$,
		such that for every $g:Fc\to d$ there are $i\in I$ and
		$\beta:c\to c_i$ for which $g$ factors as
		$Fc\xrightarrow{F\beta}Fc_i\xrightarrow{g_i}d$.
	\end{enumerate}
\end{theorem}
% segment-id: o014.li2.chapter1.adjoint-functor-theorem.pf005
\begin{proof}
	We prove only (i). Fix $c\in\Obj(\mathcal{C})$. By Lemma
	\ref{prop:adjoint-vs-comma}, $G$ has a left adjoint if and only if $(c/G)$
	has an initial object. The solution set condition says exactly that $(c/G)$
	has the weakly initial family
	$\{(d_i,c\to Gd_i):i\in I\}$, where $I$ is small. Since every initial
	object is weakly initial, the ``only if'' direction is clear. In the other
	direction, Lemma \ref{prop:jointly-weak-initial-obj} reduces the problem to
	the completeness of $(c/G)$, which is precisely Proposition
	\ref{prop:comma-cat-completeness}(ii).
\end{proof}

% segment-id: o014.li2.chapter1.adjoint-functor-theorem.e001
\begin{example}\label{eg:free-group-GAFT}
	Consider the category of groups $\cate{Grp}$ and the forgetful functor
	$U:\cate{Grp}\to\cate{Set}$. The category $\cate{Grp}$ is complete and
	$U$ preserves all small limits. To deduce from Theorem \ref{prop:GAFT}
	that $U$ has the left adjoint $\mathrm{F}:\cate{Set}\to\cate{Grp}$, which
	constructs the free group on a set, the key is to verify the solution set
	condition for every small set $c$.

	Take any $d\in\Obj(\cate{Grp})$ and map $f:c\to Ud$. Factor $f$ as
	$c\to Us\xrightarrow{U\iota}Ud$, where $s$ is the subgroup of $d$
	generated by $f(c)$ and $\iota:s\to d$ is the inclusion. Presenting $s$ by
	the generating set $f(c)$ and relations shows that the isomorphism classes
	of all such data $(s,c\to Us)$ form a small set $I$. This establishes the
	solution set condition.

	The same procedure works for other algebraic structures, such as
	$\cate{Ab}$ and $R\dcate{Mod}$, again giving left adjoints to the forgetful
	functors; see the latter part of \cite[Chapter V, \S 6]{ML98}. Likewise,
	for any commutative ring $\Bbbk$, taking the multiplicative monoid of a
	$\Bbbk$-algebra gives a forgetful functor
	$U:\Bbbk\dcate{Alg}\to\cate{Mon}$. The same technique shows that $U$ has
	a left adjoint sending a monoid $M$ to $\Bbbk[M]$; see
	\cite[Definition 5.6.1]{Li1}.

	The approach through the Adjoint Functor Theorem \ref{prop:GAFT} not only
	applies uniformly to many algebraic structures; in the free-group case it
	is also shorter than the classical construction in \cite[\S 4.8]{Li1}.
\end{example}

% segment-id: o014.li2.chapter1.adjoint-functor-theorem.p007
The Special Adjoint Functor Theorem \ref{prop:SAFT}, to be introduced shortly,
replaces the solution set condition by conditions on generators and
subobjects. We pause to define generators and cogenerators in a category;
these notions will also be used in the later section on Grothendieck
categories, among other places.

% segment-id: o014.li2.chapter1.adjoint-functor-theorem.d002
\begin{definition}\label{def:generators}
	\index{generator}
	\index{cogenerator}
	Let $\mathcal{E}$ be a category and let $\Sigma$ be a nonempty small
	subset of $\Obj(\mathcal{E})$.
% segment-id: o014.li2.chapter1.adjoint-functor-theorem.l003
	\begin{itemize}
% segment-id: o014.li2.chapter1.adjoint-functor-theorem.i006
		\item The set $\Sigma$ is called a \emph{generating family} if the
		following holds: given any pair $f,g:x\to y$ in $\mathcal{E}$, if
		$f\epsilon=g\epsilon$ for every $s\in\Sigma$ and every
		$\epsilon:s\to x$, then $f=g$. If $\{s\}$ is a generating family for
		$\mathcal{E}$, then $s$ is called a \emph{generator} of $\mathcal{E}$.
% segment-id: o014.li2.chapter1.adjoint-functor-theorem.i007
		\item The set $\Sigma$ is called a \emph{cogenerating family} if the
		following holds: given any pair $f,g:x\to y$ in $\mathcal{E}$, if
		$\delta f=\delta g$ for every $s\in\Sigma$ and every
		$\delta:y\to s$, then $f=g$. If $\{s\}$ is a cogenerating family for
		$\mathcal{E}$, then $s$ is called a \emph{cogenerator} of $\mathcal{E}$.
	\end{itemize}
\end{definition}

% segment-id: o014.li2.chapter1.adjoint-functor-theorem.p008
The two notions are plainly dual.

% segment-id: o014.li2.chapter1.adjoint-functor-theorem.le004
\begin{lemma}\label{prop:generator-via-limit}
	Let $\Sigma$ be a small subset of $\Obj(\mathcal{C})$. If every object of
	$\mathcal{C}$ can be expressed as a colimit (or limit) of elements of
	$\Sigma$, then $\Sigma$ is a generating (or cogenerating) family.
\end{lemma}
% segment-id: o014.li2.chapter1.adjoint-functor-theorem.pf006
\begin{proof}
	For colimits, a morphism $\varinjlim\alpha=:x\to y$ is uniquely determined
	by its composites with all $\alpha(i)\xrightarrow{\iota_i}\varinjlim\alpha$.
\end{proof}

% segment-id: o014.li2.chapter1.adjoint-functor-theorem.e002
\begin{example}\label{eg:generators}
	Here are some basic examples of generators and cogenerators.
% segment-id: o014.li2.chapter1.adjoint-functor-theorem.l004
	\begin{itemize}
% segment-id: o014.li2.chapter1.adjoint-functor-theorem.i008
		\item In $\cate{Set}$, the singleton $\{\mathrm{pt}\}$ is a generator:
		for maps $f,g:X\to Y$, specifying $\epsilon:\{\mathrm{pt}\}\to X$
		amounts to specifying $x\in X$, so $f\epsilon=g\epsilon$ for every
		$\epsilon$ is equivalent to pointwise equality of $f$ and $g$. The set
		$\{0,1\}$ is a cogenerator, because if $f(x)\ne g(x)$ for some $x\in X$,
		one can choose $\delta:Y\to\{0,1\}$ such that
		$\delta(f(x))\ne\delta(g(x))$.

% segment-id: o014.li2.chapter1.adjoint-functor-theorem.i009
		\item Let $\cate{CHaus}$ be the category of compact Hausdorff spaces and
		continuous maps. The singleton $\{\mathrm{pt}\}$ is again a generator,
		while $[0,1]$ is a cogenerator: for every compact Hausdorff space $Y$ and
		distinct $a,b\in Y$, Urysohn's lemma
		\cite[Theorem 6.3.1 and Corollary 7.2.6]{Xiong} gives a continuous map
		$\delta:Y\to[0,1]$ with $\delta(a)=0$ and $\delta(b)=1$.

% segment-id: o014.li2.chapter1.adjoint-functor-theorem.i010
		\item Let $R$ be a ring. In $R\dcate{Mod}$, the module $R$ is a
		generator because specifying an $R$-module homomorphism $R\to M$ amounts
		to specifying an element of $M$. Taking coproducts shows that every
		nonzero free $R$-module is also a generator.

% segment-id: o014.li2.chapter1.adjoint-functor-theorem.i011
		\item Consider the Yoneda embedding
		$h_{\mathcal{E}}:\mathcal{E}\to\mathcal{E}^\wedge$ of a small category
		$\mathcal{E}$. Then $h_{\mathcal{E}}(\Obj(\mathcal{E}))$ is a
		generating family for $\mathcal{E}^\wedge$. This follows immediately
		from the density theorem for the Yoneda embedding in Appendix \S~A.1 and
		Lemma \ref{prop:generator-via-limit}.
	\end{itemize}
\end{example}

% segment-id: o014.li2.chapter1.adjoint-functor-theorem.p009
The Special Adjoint Functor Theorem involves one more set-theoretic notion.
Recall the posets $(\mathrm{Sub}_e,\subset)$ and
$(\mathrm{Quot}_e,\twoheadleftarrow)$ introduced in Definition
\ref{def:sub-quot-obj}, where $e$ is an object of a category $\mathcal{E}$.

% segment-id: o014.li2.chapter1.adjoint-functor-theorem.d003
\begin{definition}\label{def:well-powered}
	\index{category!well-powered, well-copowered}
	A category $\mathcal{E}$ is called \emph{well-powered} (or
	\emph{well-copowered}) if $\mathrm{Sub}_e$ (or $\mathrm{Quot}_e$) is a
	small set for every $e\in\Obj(\mathcal{E})$.
\end{definition}

% segment-id: o014.li2.chapter1.adjoint-functor-theorem.p010
In this section, well-poweredness will be paired with completeness. Choose a
representative $s\hookrightarrow e$ of every element of $\mathrm{Sub}_e$.
Thus for any subset $S\subset\mathrm{Sub}_e$ it makes sense, independently of
the representatives chosen, to ask whether the fiber product
$\prod_{s\in S}(s\hookrightarrow e)$ exists. Dually, for a subset
$S\subset\mathrm{Quot}_e$, one may ask whether the fiber coproduct
$\coprod_{s\in S}(e\twoheadrightarrow s)$ exists.
% segment-id: o014.li2.chapter1.adjoint-functor-theorem.l005
\begin{itemize}
% segment-id: o014.li2.chapter1.adjoint-functor-theorem.i012
	\item If $\mathcal{E}$ is complete, then well-poweredness implies that for
	every $e$, all subsets $S\subset\mathrm{Sub}_e$ have fiber products.
% segment-id: o014.li2.chapter1.adjoint-functor-theorem.i013
	\item If $\mathcal{E}$ is cocomplete, then well-copoweredness implies that
	for every $e$, all subsets $S\subset\mathrm{Quot}_e$ have fiber coproducts.
\end{itemize}

% segment-id: o014.li2.chapter1.adjoint-functor-theorem.le005
\begin{lemma}\label{prop:cogenerator-initial-obj}
	Let $\mathcal{E}$ be complete and well-powered. If there is a cogenerating
	family $\Sigma\subset\Obj(\mathcal{E})$, then $\mathcal{E}$ has an initial
	object.
\end{lemma}
% segment-id: o014.li2.chapter1.adjoint-functor-theorem.pf007
\begin{proof}
	Take $e:=\prod_{s\in\Sigma}s$, and let $x$ be the fiber product of all
	elements of $\mathrm{Sub}_e$; only here is well-poweredness used. We show
	that $x$ is initial in $\mathcal{E}$.

	By construction, $x\hookrightarrow e$ is the infimum of the poset
	$(\mathrm{Sub}_e,\subset)$; in particular, $\mathrm{Sub}_x=\{x\}$. For
	any object $y$ and $f,g\in\Hom_{\mathcal{E}}(x,y)$, the equalizer
	$\Ker(f,g)$ must therefore be $x$ itself, so $f=g$. It remains to show that
	$\Hom_{\mathcal{E}}(x,y)$ is nonempty for every $y$.

	Form in $\mathcal{E}$ the product
	$\prod_{s\in\Sigma}s^{\Hom_{\mathcal{E}}(y,s)}$. Define
	$i:y\to\prod_{s\in\Sigma}s^{\Hom_{\mathcal{E}}(y,s)}$ so that its
	projection onto the component corresponding to $s\in\Sigma$ and
	$\delta\in\Hom_{\mathcal{E}}(y,s)$ is exactly $\delta$. Since $\Sigma$ is
	cogenerating, $i$ is easily seen to be monic. For every $s\in\Sigma$, take
	the diagonal morphism $d_s:s\to s^{\Hom_{\mathcal{E}}(y,s)}$, and form the
	pullback
% segment-id: o014.li2.chapter1.adjoint-functor-theorem.g007
	\[\begin{tikzcd}[column sep=large]
		z \arrow[r] \arrow[d, "j"'] & y \arrow[hookrightarrow, d, "i"] \arrow[phantom, ld, "\Box" description] \\
		e \arrow[r, "{\prod_{s \in \Sigma} d_s}"'] &  \prod_{s \in \Sigma} s^{\Hom_{\mathcal{E}}(y, s)}.
	\end{tikzcd}\]
	Lemma \ref{prop:mono-pullback} ensures that $j$ is also monic, giving a
	subobject of $e$. Since $x$ is the infimum of $\mathrm{Sub}_e$, there is a
	morphism $x\hookrightarrow z$. The composite $x\hookrightarrow z\to y$ is
	the required morphism.
\end{proof}

% segment-id: o014.li2.chapter1.adjoint-functor-theorem.th002
\begin{theorem}[Special Adjoint Functor Theorem]\label{prop:SAFT}
	\index{adjoint functor theorem}
	Let $\mathcal{C}$ and $\mathcal{D}$ be categories.
% segment-id: o014.li2.chapter1.adjoint-functor-theorem.l006
	\begin{enumerate}[(i)]
% segment-id: o014.li2.chapter1.adjoint-functor-theorem.i014
		\item A functor $G:\mathcal{D}\to\mathcal{C}$ has a left adjoint if it
		satisfies the following conditions.
% segment-id: o014.li2.chapter1.adjoint-functor-theorem.l007
		\begin{compactitem}
% segment-id: o014.li2.chapter1.adjoint-functor-theorem.i015
			\item $\mathcal{D}$ is complete and well-powered, and $G$ preserves
			all small limits;
% segment-id: o014.li2.chapter1.adjoint-functor-theorem.i016
			\item there is a cogenerating family
			$\Sigma\subset\Obj(\mathcal{D})$.
		\end{compactitem}
% segment-id: o014.li2.chapter1.adjoint-functor-theorem.i017
		\item A functor $F:\mathcal{C}\to\mathcal{D}$ has a right adjoint if it
		satisfies the following conditions.
% segment-id: o014.li2.chapter1.adjoint-functor-theorem.l008
		\begin{compactitem}
% segment-id: o014.li2.chapter1.adjoint-functor-theorem.i018
			\item $\mathcal{C}$ is cocomplete and well-copowered, and $F$
			preserves all small colimits;
% segment-id: o014.li2.chapter1.adjoint-functor-theorem.i019
			\item there is a generating family
			$\Sigma\subset\Obj(\mathcal{C})$.
		\end{compactitem}
	\end{enumerate}
\end{theorem}
% segment-id: o014.li2.chapter1.adjoint-functor-theorem.pf008
\begin{proof}
	By duality, it suffices to prove (i). Lemma
	\ref{prop:adjoint-vs-comma} again reduces the problem to showing that
	$(c/G)$ has an initial object for $c\in\Obj(\mathcal{C})$; the strategy is
	to apply Lemma \ref{prop:cogenerator-initial-obj}.

	Proposition \ref{prop:comma-cat-completeness}(ii) shows that $(c/G)$ is
	complete. Moreover, Definition \ref{def:generators} and the definition of
	$(c/G)$ show at once that all data $(d,c\to Gd)$ with $d\in\Sigma$ form a
	cogenerating family for $(c/G)$. Since $\Sigma$ is small, this is a small
	set.

	Let $e=(d,c\to Gd)\in\Obj((c/G))$. Subobjects are isomorphism classes of
	monomorphisms. By Proposition \ref{prop:comma-cat-completeness}, the
	projection $\Pi_{c/}:(c/G)\to\mathcal{D}$ is conservative and preserves
	monomorphisms, so it induces an order-preserving embedding
	$\mathrm{Sub}_e\hookrightarrow\mathrm{Sub}_d$. Thus $(c/G)$ is
	well-powered.
\end{proof}

% segment-id: o014.li2.chapter1.adjoint-functor-theorem.co001
\begin{corollary}\label{prop:SAFT-completeness}
	If a complete well-powered category $\mathcal{D}$ has a cogenerating family
	$\Sigma$, then $\mathcal{D}$ is cocomplete.

	Dually, if a cocomplete well-copowered category $\mathcal{C}$ has a
	generating family $\Sigma$, then $\mathcal{C}$ is complete.
\end{corollary}
% segment-id: o014.li2.chapter1.adjoint-functor-theorem.pf009
\begin{proof}
	It suffices to prove the assertion for $\mathcal{D}$. Let $I$ be a small
	category. By Example \ref{eg:comma-vs-cone}, it is enough to show that
	$(\alpha/\Delta)$ has an initial object for every
	$\alpha:I\to\mathcal{D}$. By Lemma \ref{prop:adjoint-vs-comma}, this is
	equivalent to saying that the diagonal functor $\Delta$ has a left adjoint.
	Clearly $\Delta$ preserves all small limits and colimits, so Theorem
	\ref{prop:SAFT} applies.
\end{proof}

% segment-id: o014.li2.chapter1.adjoint-functor-theorem.e003
\begin{example}
	Let $\iota:\cate{CHaus}\to\cate{Top}$ be the inclusion of compact Hausdorff
	spaces into topological spaces. Example \ref{eg:generators} showed that
	$\cate{CHaus}$ has the cogenerator $[0,1]$. By the Tychonoff theorem
	\cite[Theorem 7.7.2]{Xiong}, $\cate{CHaus}$ is readily seen to be complete,
	and $\iota$ preserves all small limits; details are left to the reader.
	Furthermore, the power set of a small set is small, so $\cate{CHaus}$ is
	well-powered. Theorem \ref{prop:SAFT}(i) therefore gives a left adjoint
	$C:\cate{Top}\to\cate{CHaus}$ to $\iota$. In point-set topology, $C(X)$
	is called the \emph{Stone--\v{C}ech compactification} of $X$; the unit gives
	a canonical morphism $\eta_X:X\to\iota C(X)$. The constructions in Lemma
	\ref{prop:cogenerator-initial-obj} and Theorem \ref{prop:SAFT} agree with
	the classical construction of the Stone--\v{C}ech compactification.
\end{example}

% segment-id: o014.li2.chapter1.adjoint-functor-theorem.p011
Adjoint functor theorems provide sufficient conditions for representability.
For representable functors, see \cite[Definition 2.5.2]{Li1} or the definition
of a representable functor in Appendix \S~A.1. The key observation is this: if
$G:\mathcal{D}\to\cate{Set}$ has a left adjoint $F$, then $G$ is
representable. Indeed, consider the singleton $\{\mathrm{pt}\}$ and put
$r_G:=F(\{\mathrm{pt}\})$. The canonical bijections
% segment-id: o014.li2.chapter1.adjoint-functor-theorem.q005
\[ \Hom_{\mathcal{D}}\left( r_G, d \right) \simeq \Hom_{\cate{Set}}\left( \{\mathrm{pt}\}, Gd \right) \simeq Gd , \quad d \in \Obj(\mathcal{D}) \]
imply $G\simeq\Hom_{\mathcal{D}}(r_G,\cdot)$.

% segment-id: o014.li2.chapter1.adjoint-functor-theorem.co002
\begin{corollary}\label{prop:cogenerator-representability}
	If a complete well-powered category $\mathcal{D}$ has a cogenerating family
	$\Sigma$, then a functor $G:\mathcal{D}\to\cate{Set}$ is representable if
	and only if it preserves all small limits.
\end{corollary}
% segment-id: o014.li2.chapter1.adjoint-functor-theorem.pf010
\begin{proof}
	The ``only if'' direction is a general property of representable functors.
	For the converse, Theorem \ref{prop:SAFT} ensures that $G$ has a left
	adjoint $F:\cate{Set}\to\mathcal{D}$; now apply the observation above.
\end{proof}

% segment-id: o014.li2.chapter1.adjoint-functor-theorem.x001
\begin{Exercises}
% segment-id: o014.li2.chapter1.adjoint-functor-theorem.ex001
	\item Consider the category of groups $\cate{Grp}$. Prove that a group
	homomorphism $\varphi:H\to G$ is a monomorphism (or epimorphism) in
	$\cate{Grp}$ if and only if its underlying map is injective (or
	surjective). Some references, including \cite{Li1}, do not adequately
	explain this fact.

% segment-id: o014.li2.chapter1.adjoint-functor-theorem.hn001
	\begin{hint}
		It suffices to prove the ``only if'' direction. For monomorphisms,
		consider homomorphisms from $\Ker\varphi$ to $H$. For epimorphisms, one
		must show that if $\varphi(H)\ne G$, then there are distinct group
		homomorphisms $f,g:G\rightrightarrows K$ whose restrictions to
		$\varphi(H)$ agree. One approach is to take $K$ to be the amalgamated
		free product defined by two copies of the embedding
		$\varphi(H)\hookrightarrow G$ in \cite[Definition 4.8.5]{Li1}, and to
		take $f$ and $g$ as the embeddings of $G$ into the two copies. Uniqueness
		of reduced expressions \cite[Lemma 4.8.7]{Li1} implies $f\ne g$.
	\end{hint}

% segment-id: o014.li2.chapter1.adjoint-functor-theorem.ex002
	\item Show that in the category $\cate{Top}$, the image of $f:X\to Y$ is
	$f(X)$ with the subspace topology from $Y$, while its coimage is $f(X)$
	with the quotient topology from $X$. Give the corresponding topological
	interpretations of strict monomorphisms and strict epimorphisms.

% segment-id: o014.li2.chapter1.adjoint-functor-theorem.ex003
	\item Let $\mathcal{C}$ have finite limits and finite colimits. Prove that
	the following statements about a morphism $f:X\to Y$ are equivalent.
% segment-id: o014.li2.chapter1.adjoint-functor-theorem.l009
	\begin{enumerate}[(i)]
% segment-id: o014.li2.chapter1.adjoint-functor-theorem.i020
		\item The morphism $f$ is a strict epimorphism.
% segment-id: o014.li2.chapter1.adjoint-functor-theorem.i021
		\item The canonical morphism $\Coim(f)\to Y$ is an isomorphism.
% segment-id: o014.li2.chapter1.adjoint-functor-theorem.i022
		\item $X\dtimes{Y}X\rightrightarrows X\xrightarrow{f}Y$ is a coequalizer
		diagram.
% segment-id: o014.li2.chapter1.adjoint-functor-theorem.i023
		\item The morphism $f:X\to Y$ is the coequalizer of some pair
		$a,b:W\rightrightarrows X$.
	\end{enumerate}
% segment-id: o014.li2.chapter1.adjoint-functor-theorem.hn002
	\begin{hint}
		Recall that being epic is equivalent to $\Image(f)\rightiso Y$ (Lemma
		\ref{prop:epi-image}), and use Proposition
		\ref{prop:epi-mono-isomorphism}. It is easy to prove
		(i) $\iff$ (ii) $\iff$ (iii) $\implies$ (iv). For
		(iv) $\implies$ (ii), consider the solid part of the commutative diagram
% segment-id: o014.li2.chapter1.adjoint-functor-theorem.g008
		\[\begin{tikzcd}
			W \arrow[shift left, r, "a"] \arrow[shift right, r, "b"'] \arrow[d, "{(a, b)}"'] & X \arrow[d, "\identity"] \arrow[twoheadrightarrow, r, "f"] & Y \arrow[dashed, d] \\
			X \dtimes{Y} X \arrow[shift left, r] \arrow[shift right, r] & X \arrow[twoheadrightarrow, r] & \Coim(f)
		\end{tikzcd}\]
		The universal property of the coequalizer gives the dashed arrow and
		shows that it and $\Coim(f)\to Y$ are inverse to each other.
	\end{hint}

% segment-id: o014.li2.chapter1.adjoint-functor-theorem.ex004
	\item Let $\mathcal{C}$ have finite limits and finite colimits. Prove that
	if a morphism $f$ has a right inverse (or left inverse), then $f$ is a
	strict epimorphism (or strict monomorphism).
% segment-id: o014.li2.chapter1.adjoint-functor-theorem.hn003
	\begin{hint}
		If $fg=\identity_Y$, then $f$ is the coequalizer of $gf$ and
		$\identity_X$.
	\end{hint}

% segment-id: o014.li2.chapter1.adjoint-functor-theorem.ex005
	\item Prove that an additive-category structure on any category, if it
	exists, is unique.
% segment-id: o014.li2.chapter1.adjoint-functor-theorem.hn004
	\begin{hint}
		Apply Proposition \ref{prop:additive-cat-addition}.
	\end{hint}

% segment-id: o014.li2.chapter1.adjoint-functor-theorem.ex006
	\item Compare Example \ref{eg:Grp-Set-create}. Let $U$ be the forgetful
	functor from compact Hausdorff spaces $\cate{CHaus}$ to sets $\cate{Set}$.
	Prove that $U$ creates all small limits. Show that it does not create all
	small colimits.

% segment-id: o014.li2.chapter1.adjoint-functor-theorem.ex007
	\item Let $\mathcal{I}$ be the category of all finite sets and surjective
	maps. Prove that $\mathcal{I}^{\opp}$ is not filtered in the sense of
	\S\ref{sec:filtered-indlim}.

% segment-id: o014.li2.chapter1.adjoint-functor-theorem.ex008
	\item In Definition \ref{def:Kan-extension}, prove that every morphism
	$F\to F'$ induces canonical morphisms $\Lan_KF\to\Lan_KF'$ and
	$\Ran_KF\to\Ran_KF'$, provided these Kan extensions exist. Characterize
	the induced morphisms.

% segment-id: o014.li2.chapter1.adjoint-functor-theorem.ex009
	\item Let $K$ and $F$ be as in \S\ref{sec:Kan-extension}. Prove that if
	$M:\mathcal{E}\to\mathcal{F}$ has a left adjoint (or right adjoint), then
	$M$ preserves $\Ran_KF$ (or $\Lan_KF$).

% segment-id: o014.li2.chapter1.adjoint-functor-theorem.hn005
	\begin{hint}
		By duality, suppose that $M$ has a right adjoint $N$, with unit
		$\eta':\identity\to NM$ and counit $\varepsilon':MN\to\identity$. For
		any functor $L:\mathcal{D}\to\mathcal{F}$, there are canonical
		bijections
% segment-id: o014.li2.chapter1.adjoint-functor-theorem.q006
		\begin{multline*}
			\Hom_{\mathcal{F}^{\mathcal{D}}} \left( M (\Lan_K F) , L \right) \simeq \Hom_{\mathcal{E}^{\mathcal{D}}} \left( \Lan_K F, N L \right) \\
			\stackrel{\text{\eqref{eqn:Kan-univ}}}{\simeq} \Hom_{\mathcal{E}^{\mathcal{C}}} \left( F, NLK \right) \simeq \Hom_{\mathcal{F}^{\mathcal{C}}} \left( MF, LK \right).
		\end{multline*}

		Using the general properties of $\eta'$ and $\varepsilon'$, verify that
		when $L=M(\Lan_KF)$ this bijection sends $\identity$ to $M\eta$, where
		$\eta$ is the morphism belonging to $\Lan_KF$.

		For general $L$, use the familiar technique from the proof of the Yoneda
		lemma to verify that an arbitrary morphism
		\[ \varphi:M(\Lan_KF)\to L \]
		corresponds to the composite
% segment-id: o014.li2.chapter1.adjoint-functor-theorem.q007
		\[
		\begin{aligned}
			(\varphi K)(M\eta):\quad MF
			&\xrightarrow{M\eta} M(\Lan_K F)K, \\
			M(\Lan_K F)K &\xrightarrow{\varphi K} LK .
		\end{aligned}
		\]
	\end{hint}

% segment-id: o014.li2.chapter1.adjoint-functor-theorem.ex010
	\item Consider the diagram of functors
% segment-id: o014.li2.chapter1.adjoint-functor-theorem.g009
	\[\begin{tikzcd}[column sep=small, row sep=small]
		& \mathcal{D} \arrow[rd, "E"] & \\
		\mathcal{C} \arrow[ru, "F"] & & \mathcal{D}' \arrow[ll, "G"]
	\end{tikzcd}\]
	where $EF$ is left adjoint to $G$, with counit
	$\varepsilon:EFG\to\identity_{\mathcal{D}'}$. Prove that if, for every
	category $\mathcal{X}$, the functor
	$F^*:\mathcal{X}^{\mathcal{D}}\to\mathcal{X}^{\mathcal{C}}$ is fully
	faithful, then $E$ can also be realized as a left adjoint to $FG$, with the
	same counit $\varepsilon$.

% segment-id: o014.li2.chapter1.adjoint-functor-theorem.hn006
	\begin{hint}
		The following argument is from \cite[I.1.3.1 Lemma]{GZ67}. Let
		$\eta:\identity_{\mathcal{C}}\to GEF$ be the unit of $(EF,G)$. Taking
		$\mathcal{X}=\mathcal{D}$ gives
		$\eta':\identity_{\mathcal{D}}\to FGE$ such that $\eta'F=F\eta$. Check
		the triangle identities for $(E,FG)$ with $\eta'$ and $\varepsilon$.
		First,
		$(FG\varepsilon)(\eta'FG)=(FG\varepsilon)(F\eta G)=\identity_{FG}$.
		On the other hand, taking $\mathcal{X}=\mathcal{D}'$ shows that
		$(\varepsilon E)(E\eta')=\identity_E$ is equivalent to
		$(\varepsilon EF)(E\eta'F)=\identity_{EF}$; but
		$(\varepsilon EF)(E\eta'F)=(\varepsilon EF)(EF\eta)=\identity_{EF}$.
	\end{hint}

% segment-id: o014.li2.chapter1.adjoint-functor-theorem.ex011
	\item Consider an adjoint pair
% segment-id: o014.li2.chapter1.adjoint-functor-theorem.g010
	$\begin{tikzcd}
		F: \mathcal{C} \arrow[r, shift left] & \mathcal{D} \arrow[l, shift left] : G
	\end{tikzcd}$,
	and put
% segment-id: o014.li2.chapter1.adjoint-functor-theorem.q008
	\[ S := \left\{ s \in \Mor(\mathcal{C}): Fs \;\text{is invertible} \right\}; \]
	this factors $F$ uniquely as
	$\mathcal{C}\xrightarrow{Q}\mathcal{C}[S^{-1}]\xrightarrow{H}\mathcal{D}$.
	Prove that the following statements are equivalent.
% segment-id: o014.li2.chapter1.adjoint-functor-theorem.l010
	\begin{enumerate}[(i)]
% segment-id: o014.li2.chapter1.adjoint-functor-theorem.i024
		\item The functor $G$ is fully faithful.
% segment-id: o014.li2.chapter1.adjoint-functor-theorem.i025
		\item The counit $\varepsilon:FG\to\identity_{\mathcal{D}}$ is an
		isomorphism.
% segment-id: o014.li2.chapter1.adjoint-functor-theorem.i026
		\item The functor $H$ is an equivalence of categories, with $QG$ as a
		quasi-inverse.
% segment-id: o014.li2.chapter1.adjoint-functor-theorem.i027
		\item For every category $\mathcal{X}$, the functor
		$F^*:\mathcal{X}^{\mathcal{D}}\to\mathcal{X}^{\mathcal{C}}$ is fully
		faithful.
	\end{enumerate}
	Also state the dual version of this result.

% segment-id: o014.li2.chapter1.adjoint-functor-theorem.hn007
	\begin{hint}
		The following argument is from \cite[I.1.3 Proposition]{GZ67}. First,
		(i) $\iff$ (ii) is fairly easy; see the hint in
		\cite[Chapter 2, Exercise 8]{Li1}.

		For (ii) $\implies$ (iii), one must prove
		$QGH\simeq\identity_{\mathcal{C}[S^{-1}]}$. Proposition
		\ref{prop:localization-extra-property} reduces this to
		$QGHQ\simeq Q$. Consider the unit
		$\eta:\identity_{\mathcal{C}}\to GF$. Since
		$(\varepsilon F)(F\eta)=\identity_F$, the morphism $F\eta$ is an
		isomorphism, hence $Q\eta:Q\rightiso QGF=QGHQ$.

		For (iii) $\implies$ (iv), the fact that $H$ is an equivalence implies
		that $H^*$ is fully faithful, while Proposition
		\ref{prop:localization-extra-property} applies to $Q^*$.

		For (iv) $\implies$ (ii), put $E=\identity_{\mathcal{D}}$ in the
		preceding exercise. Then $FG$ is a right adjoint to
		$\identity_{\mathcal{D}}$, with counit $\varepsilon$. Since right
		adjoints to $\identity_{\mathcal{D}}$ are unique up to isomorphism,
		$\varepsilon$ is an isomorphism.
	\end{hint}

% segment-id: o014.li2.chapter1.adjoint-functor-theorem.ex012
	\item Let $S$ be a left multiplicative system in $\mathcal{C}$. Prove that
	every commutative diagram in $\mathcal{C}[S^{-1}]^l$ of the form
% segment-id: o014.li2.chapter1.adjoint-functor-theorem.g011
	\[\begin{tikzcd}
		Q^l X \arrow[r, "{Q^l f}"] \arrow[d, "\alpha"'] & Q^l X' \arrow[d, "\beta"] \\
		Q^l Y \arrow[r, "{Q^l g}"'] & Q^l Y'
	\end{tikzcd}\]
	\begin{center}
		arises from a commutative diagram in $\mathcal{C}$ of the form
	\end{center}
% segment-id: o014.li2.chapter1.adjoint-functor-theorem.g012
	\[\begin{tikzcd}
		X \arrow[r, "f"] & X' \\
		U \arrow[r] \arrow[rightarrowtail, u] \arrow[d] & U' \arrow[rightarrowtail, u] \arrow[d] \\
		Y \arrow[r, "g"'] & Y'
	\end{tikzcd}\]
	where $f,g\in\Mor(\mathcal{C})$ are given and, as usual,
	$\rightarrowtail$ denotes a morphism in $S$. State the corresponding
	version for a right multiplicative system.

% segment-id: o014.li2.chapter1.adjoint-functor-theorem.hn008
	\begin{hint}
		Represent $\alpha$ and $\beta$ by $(U^\dagger;s,a)$ and $(U';t,b)$,
		respectively. Apply (S3) to obtain $W,h,r$ and a commutative diagram in
		$\mathcal{C}$
% segment-id: o014.li2.chapter1.adjoint-functor-theorem.g013
		\[\begin{tikzcd}
			W \arrow[r, "h"] \arrow[rightarrowtail, d, "r"'] & U' \arrow[rightarrowtail, d, "{t}"] \\
			U^\dagger \arrow[r, "{fs}"'] & X' .
		\end{tikzcd}\]
		This shows that the upper square in the following diagram commutes:
% segment-id: o014.li2.chapter1.adjoint-functor-theorem.g014
		\[\begin{tikzcd}
			X \arrow[r, "f"] & X' \\
			W \arrow[rightarrowtail, u, "{sr}"] \arrow[r, "h"] \arrow[d, "{ar}"'] & U' \arrow[rightarrowtail, u, "{t}"'] \arrow[d, "{b}"] \\
			Y \arrow[r, "g"'] & Y'
		\end{tikzcd}\]
		For the lower square, first show that it commutes after applying $Q^l$;
		then use Corollary \ref{prop:fraction-zero} to choose
		$U\rightarrowtail W$ and adjust it to the required commutative diagram.
	\end{hint}

% segment-id: o014.li2.chapter1.adjoint-functor-theorem.ex013
	\item Verify the details of central localization in Example
	\ref{eg:central-localization}. Prove that the category
	$\mathcal{C}[S^{-1}]$ defined there and the localization supplied by
	Theorem \ref{prop:localization-additivity} are equivalent as
	$Z(\mathcal{C})[S^{-1}]$-linear categories.

% segment-id: o014.li2.chapter1.adjoint-functor-theorem.ex014
	\item Complete the details in the proof of Proposition
	\ref{prop:comma-cat-completeness}(i).

% segment-id: o014.li2.chapter1.adjoint-functor-theorem.ex015
	\item Use the Adjoint Functor Theorem \ref{prop:GAFT} to construct the free
	product $G\star H$ of any two groups $G$ and $H$; see
	\cite[Definition 4.8.9]{Li1}.

% segment-id: o014.li2.chapter1.adjoint-functor-theorem.ex016
	\item Let $\mathrm{F}(X)$ be the free group on a set $X$. Prove that the
	canonical map $\iota:X\to U(\mathrm{F}(X))$ is injective, where $U$ is the
	forgetful functor from groups to sets.
% segment-id: o014.li2.chapter1.adjoint-functor-theorem.hn009
	\begin{hint}
		For distinct $x,y\in X$, there is a group $H$ and a map of sets
		$f:X\to U(H)$ such that $f(x)\ne f(y)$.
	\end{hint}

% segment-id: o014.li2.chapter1.adjoint-functor-theorem.ex017
	\item (P.\ Freyd) Let $\mathcal{D}$ be complete and let
	$G:\mathcal{D}\to\cate{Set}$ preserve all small limits and satisfy the
	following solution set condition: there is a small subset
	$\Gamma\subset\Obj(\mathcal{D})$ such that for every
	$d\in\Obj(\mathcal{D})$ and $p\in Gd$, there are $x\in\Gamma$, $q\in Gx$,
	and $f:x\to d$ with $(Gf)(q)=p$. Prove that $G$ is representable.

% segment-id: o014.li2.chapter1.adjoint-functor-theorem.hn010
	\begin{hint}
		Write an object of $\left(\{\mathrm{pt}\}/G\right)$ as $(d,p)$, where
		$d\in\Obj(\mathcal{D})$ and $p\in Gd$. The condition says exactly that
		$\Gamma^\natural:=\{(x,q):x\in\Gamma,\;q\in Gx\}$ is a weakly initial
		family for $\left(\{\mathrm{pt}\}/G\right)$. Show that this category
		has an initial object $(r_G,u_G)$. Following the argument for
		\eqref{eqn:adjoint-vs-comma-aux}, show that for every $d$ there is a
		bijection
% segment-id: o014.li2.chapter1.adjoint-functor-theorem.g015
		\[\begin{tikzcd}[row sep=tiny, column sep=small]
			\Hom_{\mathcal{D}}(r_G, d) \arrow[r] & \Hom_{\cate{Set}}\left( \{\mathrm{pt}\}, Gd \right) \arrow[r, "\sim"] & Gd \\
			\alpha \arrow[mapsto, rr] & & (G\alpha)(u_G) .
		\end{tikzcd}\]
	\end{hint}
\end{Exercises}
